% ID: SEL-D0074-OB-E02
% Tanda S074.
\textbf{a)} La derivada $f'(x)=3x(x-2)$ es positiva fuera de $[0,2]$ y negativa en $(0,2)$. Por tanto, $f$ crece en $(-\infty,0)$ y $(2,\infty)$, y decrece en $(0,2)$.

Como $f''(x)=6(x-1)$, la gráfica se curva hacia abajo en $(-\infty,1)$ y hacia arriba en $(1,\infty)$. El cambio de signo en $1$ da el punto de inflexión $(1,2)$.

\textbf{b)} En $x=0$ la derivada pasa de positiva a negativa y en $x=2$ pasa de negativa a positiva. Sustituyendo en la función, el máximo relativo es $\boxed{(0,4)}$ y el mínimo relativo es $\boxed{(2,0)}$.

\textbf{c)} Como $f(-1)=0$ y $f'(-1)=9$, la tangente cumple $y-0=9(x+1)$, es decir, $\boxed{y=9x+9}$.
